\documentclass[10pt,a4paper]{amsart}
\usepackage[margin=19mm]{geometry}
\usepackage{amsmath,amssymb,mathtools}
\usepackage[scr=rsfs,cal=euler]{mathalfa}
\usepackage{xfrac}
\usepackage[unicode,hidelinks,bookmarks=false]{hyperref}
\newcommand{\ie}{i.e.\ }
\newcommand{\cf}{cf.\ }
\newcommand{\resp}{respectively\ }
\newcommand{\Subsection}{\subsection}
\providecommand{\nobreakdash}{\nobreak\hbox{-}}
\setlength{\emergencystretch}{2em}
\multlinegap0pt
\usepackage{amssymb}
\usepackage{tikz}
\usepackage[shortlabels]{enumitem}
\usepackage{float}
\usepackage{mathtools}
\usepackage{cancel}
\newtheorem{lemma}{Lemma}
\newcommand{\ang}[1]{\langle #1\rangle}
\title{Regularization by noise for Gevrey well-posedeness of a weakly hyperbolic operator}
\author{Enrico Bernardi, Alberto Lanconelli}
\date{Source: arXiv:2603.05050v1 (CC BY 4.0)}
\begin{document}
\maketitle

\noindent\emph{Source and licence.} Regularization by noise for Gevrey well-posedeness of a weakly hyperbolic operator. Version arXiv:2603.05050v1 (CC BY 4.0), distributed by arXiv under the Creative Commons Attribution 4.0 International licence, http://creativecommons.org/licenses/by/4.0/, as recorded in the arXiv licence metadata for this exact version and shown on the abstract page https://arxiv.org/abs/2603.05050v1. Full licence notice: \texttt{licenses/arxiv-2603.05050-CC-BY-4.0.txt}.\par\smallskip

\noindent\textit{Editorial scope.} Counted unit: the article's lemma lem:qF (source0.tex lines 744--749). The article's complete original proof of it is delivered with it (source0.tex lines 751--792, one \texttt{proof} environment of 42 lines). No mathematics is written, repaired or re-derived: the statement and the proof are the article's own lines, in the article's own notation, and no retained line is substituted or re-flowed.  Retained uncounted as the source-local context the proof needs: lines 703--711; lines 713--715; lines 717--719; lines 721--730.  Nothing was omitted from any retained span, so the package carries no omission record.  No retained line contains a citation, so the wrapper carries no bibliography and no bibliography entry.  No retained line contains a figure environment, an image-inclusion command or drawing code, so no image is delivered and the package declares no asset dependency.  Editorial formatting only, in addition: an AMS \texttt{amsart} preamble that restates the article's own declarations and macros used by the retained lines, namely the environments \texttt{lemma} on separate counters, together with the standard packages those lines need.  The retained environments print in this document as Lemma 1 in order of appearance, whereas the article prints the same items with the numbering of its own full document.  The complete native source of the article as distributed in the arXiv e-print for this exact version is bundled unchanged as \texttt{sources/195781/source0.tex}.
\begin{align}
q_0(0) &= \frac12 m_{10}-\frac{1}{2\eta}m_{30}, \label{eq:q0}\\
q_+(0) &= -\frac{\lambda_-}{2\Delta}m_{10}
         +\frac{2}{\Delta}m_{20}
         -\frac{\lambda_-}{2\eta\Delta}m_{30},\label{eq:qp}\\
q_-(0) &= \frac{\lambda_+}{2\Delta}m_{10}
         -\frac{2}{\Delta}m_{20}
         +\frac{\lambda_+}{2\eta\Delta}m_{30}. \label{eq:qm}
\end{align}

\begin{equation}\label{eq:F0-trivial}
m_{10}\le F(0;\xi),\qquad m_{30}\le \ang{\xi}\,F(0;\xi).
\end{equation}

\begin{equation}\label{eq:m20-bound}
|m_{20}|\le \frac{m_{10}+m_{30}}{2}.
\end{equation}

\begin{lemma}[Large-$|\xi|$ coefficient bounds]\label{lem:coef}
There exist $\xi_0\ge1$ and $C>0$ (depending only on $\sigma$) such that for all $|\xi|\ge\xi_0$,
\begin{equation}\label{eq:coef}
\frac1{|\Delta|}\le\frac{C}{\xi^2},\qquad
\Bigl|\frac{\lambda_-}{\Delta}\Bigr|\le C,\qquad
\Bigl|\frac{\lambda_-}{\xi\Delta}\Bigr|\le\frac{C}{|\xi|},\qquad
\Bigl|\frac{\lambda_+}{\Delta}\Bigr|\le\frac{C}{|\xi|^3},\qquad
\Bigl|\frac{\lambda_+}{\xi\Delta}\Bigr|\le\frac{C}{|\xi|^4}.
\end{equation}
\end{lemma}

\begin{lemma}[Control by $F(0;\xi)$ of eigenvector coefficients]\label{lem:qF}
There exist $\xi_0\ge1$ and $C>0$ such that for all $|\xi|\ge\xi_0$,
\begin{equation}\label{eq:qF}
|q_0(0)|+|q_+(0)|+|q_-(0)|\le C\,F(0;\xi).
\end{equation}
\end{lemma}

\begin{proof}
\emph{Step 1: bound $q_0(0)$.}
By \eqref{eq:q0} and \eqref{eq:F0-trivial},
\[
|q_0(0)|\le \tfrac12 m_{10}+\tfrac1{2|\xi|}m_{30}
\le \tfrac12F(0;\xi)+\tfrac{\ang{\xi}}{2|\xi|}\frac{m_{30}}{\ang{\xi}}
\le C F(0;\xi),
\]
since $\ang{\xi}/|\xi|\le \sqrt{2}$ for $|\xi|\ge 1$.

\medskip
\noindent
\emph{Step 2: bound $q_+(0)$.}
Using \eqref{eq:qp} and Lemma~\ref{lem:coef},
\[
|q_+(0)|\le C m_{10}+\frac{C}{\xi^2}|m_{20}|+\frac{C}{|\xi|}m_{30}.
\]
We bound each term by $F(0;\xi)$. First, $m_{10}\le F(0;\xi)$.
Next, by \eqref{eq:m20-bound} and \eqref{eq:F0-trivial},
\[
\frac{1}{\xi^2}|m_{20}|\le \frac{1}{2\xi^2}(m_{10}+m_{30})
\le \frac{C}{\xi^2}F(0;\xi)+\frac{C}{\xi^2}\ang{\xi}F(0;\xi)
\le \frac{C}{|\xi|}F(0;\xi)\le C F(0;\xi),
\]
for $|\xi|\ge 1$.
Finally,
\[
\frac{1}{|\xi|}m_{30}\le \frac{\ang{\xi}}{|\xi|}\frac{m_{30}}{\ang{\xi}}
\le C F(0;\xi).
\]
Hence $|q_+(0)|\le C F(0;\xi)$.

\medskip
\noindent
\emph{Step 3: bound $q_-(0)$.}
Using \eqref{eq:qm} and Lemma~\ref{lem:coef},
\[
|q_-(0)|\le \frac{C}{|\xi|^3}m_{10}+\frac{C}{\xi^2}|m_{20}|+\frac{C}{|\xi|^4}m_{30}
\le C F(0;\xi),
\]
using again \eqref{eq:m20-bound} and \eqref{eq:F0-trivial} and $|\xi|\ge 1$. Summing the three bounds yields \eqref{eq:qF}.
\end{proof}

\end{document}
